\section{The prototype: what exists and what it shows}\label{sec:status}

\VG{} is being built by AI engineering agents in phases, each closed by a gate: a script,
\kw{make gate-gN}, that exits $0$ only when every criterion holds and writes a report, which the
human owner must approve. The agents may not mark a gate passed. This section reports the state
at commit \texttt{fe499c8} (30 October 2026).

\subsection{What is implemented}

\begin{itemize}
\item \textbf{vg-core} (8,399 lines of Rust, plus 660 lines of integration tests): the data
  model; the VG-Sym lexer, parser and printer; lowering with the structural validator
  (arities and modes from the opcode table, index ranges, acyclicity, single assignment,
  define-before-use); the canonicalizer and the body, contract and node hashes; the binary
  codec with its 128-byte header and section directory; a disassembler; and a seeded program
  generator for property tests.
\item \textbf{vg-store} (333 lines): a content-addressed object store, a SQLite index and the
  lock file that pins label resolutions.
\item \textbf{vg-cli} (196 lines): the commands \kw{vg asm}, \kw{vg disasm} and \kw{vg hash}.
\end{itemize}
The other nine crates (types, interpreter, SMT, code generation, runtime, decoding, telemetry,
registry, lifter) are empty skeletons. There is therefore no type checker, no interpreter, no
verifier, no code generator, no agent loop and no benchmark. Every statement in this paper about
verification is a statement about the specification.

\subsection{Evidence for the core representation}

We ran \kw{cargo test --workspace --locked} at \texttt{fe499c8}: 26 tests pass and none fail.
\Cref{tab:g1} summarizes the evidence of the gate~G1 report, generated by \kw{make gate-g1}
for commit \texttt{bf3935d}, whose gate inputs are identical to those of \texttt{fe499c8}. The
report states that every G1 criterion holds; the owner has not approved it.

\begin{table}[t]
\centering
\small
\begin{tabularx}{\linewidth}{@{}l>{\raggedright\arraybackslash}X@{}}
\toprule
Law or criterion & Evidence \\
\midrule
Parse after print is the identity & 206 VG-Sym files \\
Round trips and idempotence & 205 files, 386 nodes: canonicalization idempotent,
  $\mathrm{decode}\circ\mathrm{encode}=\mathrm{id}$, $\mathrm{asm}\circ\mathrm{disasm}=\mathrm{id}$,
  $\mathrm{disasm}\circ\mathrm{asm}$ a fixed point; 7 golden binary files reproduced byte for
  byte \\
Equivalence moves keep the hash & 56 alpha-equivalence families (168 hand-written files) and 560
  seeded random variants (seed \texttt{0x5eed20261030}): all hash identically \\
Inequivalent programs differ & 16 distinct families, 18 pairs: all hash differently \\
Hashes follow unraveling & 400 generated bodies of 3 to 30 ops (seed \texttt{0xd0c20261030}):
  400 variants hash and unravel equally; of the near-misses, 163 unravel differently and hash
  differently, 559 unravel equally and hash equally \\
Parser agrees with the reference grammar & 245 files (226 accepted, 19 rejected), same verdict
  and token texts as the Lark grammar \\
Fuzzing without a crash & 4,184\,s per target: parser 13,958,693 runs, codec 56,196,128, canonicalizer
  545,030 \\
Generated programs & 600 programs (888 nodes) assembled, canonical and decoded back \\
\bottomrule
\end{tabularx}
\caption{Evidence of the G1 report for the core representation (not approved by the owner at
the time of writing).}
\label{tab:g1}
\end{table}

These are tests: they confirm the laws of \cref{sec:identity} on the inputs they cover. The
``hashes follow unraveling'' test is the empirical counterpart of \cref{thm:hash-char}: it
compares the hash with a slow reference unraveling on random bodies.

\subsection{An independent review, read through the theorems}

A separate AI session reviewed gate~G1 adversarially before approval. It reran
\kw{make gate-g1} (exit $0$), checked 23 pairs of programs with \kw{vg hash} (all behaved as the
specification requires), and wrote eight Rust tests, five of which failed. Its findings about
the core representation are:
\begin{itemize}
\item[B1] Type member counts are written into one byte, so the two different tuple types of
  \cref{rem:capacity} serialize identically.
\item[B2] Results beyond index 255 of an op alias earlier results: \kw{RETURN(\%256)} of a call
  with 257 results returns result $0$, and hashes like it.
\item[B3] Such wide programs hash but cannot be stored: the assembler writes bytes that its own
  decoder rejects.
\item[B4] The merge step merges effect ops: two clock reads on one token become one.
\item[B5] The decoder accepts a node in which one effect token feeds two effect ops.
\end{itemize}
It also found three ways to make the gate script report criteria met while they are not (B6, of
medium severity: a cargo
configuration outside the gate's inputs that hides failing tests; B7 and B8, of low severity:
ignored files invisible to its clean-tree check, and no minimum on the number of distinct
families), and recommended not approving
G1 until B1 to B3 are fixed, since the format freezes when G1 passes.

Read through \cref{sec:identity}, B1 to B3 are exactly violations of \cref{cond:capacity}, the
side condition of \cref{lem:prefixfree} and hence of \cref{thm:hash-char}; and B4 and B5 are
exactly the failure of \ref{wf:eff} on a path into the canonicalizer, the hypothesis of
\cref{lem:eff-never-merged}. The theorems do not find such bugs by themselves; the review found
them by testing. What the theorems add is a precise account of why exactly these inputs break the
guarantee and which checks restore it: bound every count by its field, or widen the field; check
the linearity of effect tokens in the decoder as well as in the text lowering.

We reproduced B1 with the prototype's toolchain. For identity functions on the two tuple types
of \cref{rem:capacity}, \kw{vg hash} prints the same node hash, \texttt{\#e271d220c54d}, and the
same body and contract hashes, while the analogous pair with two and one inner members gets two
different node hashes; \kw{vg asm} rejects the wide program with
\texttt{VG-S0202 unknown type tag 0x00}, as B3 states (\cref{app:repro}).

\subsection{Independent checks of this paper}

The project's rules ask for independent checks, so we sent the proofs to models of other
families, each without tools and with the relevant specification text in the prompt. Gemini 3.8
Flash, through Google's Antigravity command-line agent, received four briefs.
\begin{itemize}
\item The soft-knee proof: it reached the same port-by-port conclusions as
  \cref{thm:softknee,prop:draft}, and contributed the observation on signed zeros in
  \cref{rem:contract}. One decimal expansion it gave for the binary32 value of $3\cdot10^{38}$ was
  wrong (its hexadecimal value was right); we use the exact value $14791142\cdot2^{104}$.
\item The identity theorems: it gave the same induction as \cref{thm:hash-char} and
  \cref{lem:eff-never-merged}, the same capacity counterexamples, and pointed out that the
  specification's ``exactly when'' statement about equivalence moves also needs the
  identifications we list after \cref{cor:canon}. We rejected two of its findings because
  they contradicted specification text it had not been given: a commutativity problem for
  \op{FMIN} and \op{FMAX} that holds for some hardware minimum instructions but not for IEEE
  754-2019 \emph{minimum}, and a dangling reference from an invariant that the specification
  drops together with its ops.
\item Soundness: it agreed with \cref{lem:gen,lem:weak,thm:widths} and with the fold rule, and
  independently found that facts of \kw{spec:} applications must be guarded
  (\cref{prop:foldspec}). It also found that the rules the specification gives for two
  operations planned after the MVP, \op{SCAN\_PREFIX} and \op{MATCH\_DISJOINT}, have no
  precondition goals for the operator or the arms (\cref{sec:discussion}). We rejected two
  findings: missing obligations for \op{FCOPYSIGN} and \op{FCLAMP}, which the opcode table does
  state but its prompt did not include, and a possible negative index, which the unsigned index
  type of \op{INDEX} excludes.
\item Related work and threats to validity, which shaped \cref{sec:related,sec:discussion}. No
  reference was taken from the model: references with a DOI were checked against Crossref and
  preprints against the arXiv API, and for the few others we give only details we are certain
  of.
\end{itemize}
Two further Gemini briefs reviewed the finished draft. A referee brief on the formal sections
found that \cref{tab:rules} did not tie the result of a call, map, generate or fold to the fact
about it, an omission we corrected by naming that result $p_i$ in the facts; that the
evaluation of effectful calls and guards on the effect chain needed the precise statement now in
\cref{def:eval}; and two notational slips in \cref{sec:case}. A fact-checking brief compared the
non-formal sections with the gate report, the progress log, the review and the enumeration
output, and found four overclaims in an earlier abstract, introduction and acknowledgment (a gate
described as passed although it awaits approval, evaluations counted as inputs, the extent of
the GPT-OSS reviews, and the severity grouping of the review's findings), all corrected.

GPT-OSS 120B received the soundness brief and the referee brief. The first stopped with a
capacity error after a partial answer whose claims were either consistent with ours or mistaken
(it called the fresh-element rule unsound for relational properties, which confuses completeness
with soundness). The second reported five numerical errors, all of which we checked and found
to be its own: for example, it gave the spacing of binary32 values in $[2^{127},2^{128})$ for
that of $[2^{126},2^{127})$, whose largest element is indeed $2^{127}-2^{103}$. Gemini 4 Argon was
unavailable. None of these checks is a proof; they are reviews, and every claim we adopted from
them was rechecked by hand.
